% Created 2026-04-01 Wed 23:59
% Intended LaTeX compiler: pdflatex
\documentclass[11pt, letterpaper]{article}

\usepackage[utf8]{inputenc}
\usepackage[T1]{fontenc}
\usepackage{graphicx}
\usepackage{longtable}
\usepackage{wrapfig}
\usepackage{rotating}
\usepackage[normalem]{ulem}
\usepackage{amsmath}
\usepackage{amssymb}
\usepackage{capt-of}
\usepackage{hyperref}
\usepackage[margin=2.5cm]{geometry}      % Márgenes decentes
\usepackage[utf8]{inputenc}
\usepackage{palatino}                   % Tipografía elegante
\usepackage{xcolor}                     % Colores personalizados
\author{Equipo Alpine White}
\date{\textit{{[}2026-03-21 Sat]}}
\title{Bitácora Semanal 7: Administración de Sistemas Unix}
\hypersetup{
 pdfauthor={Equipo Alpine White},
 pdftitle={Bitácora Semanal 7: Administración de Sistemas Unix},
 pdfkeywords={},
 pdfsubject={},
 pdfcreator={},
 pdflang={English}}
\begin{document}

\maketitle
\setcounter{tocdepth}{2}
\tableofcontents

\section{Información General}
\label{sec:orgd3fad95}
\begin{itemize}
\item \textbf{Semana:} Del 16 de Marzo 2026 al 20 de Marzo 2026
\item \textbf{Equipo :}
\begin{itemize}
\item Arreguín Salgado Gael Emiliano
\item Gramer Muñoz Omar Fernando
\item López Pérez Mariana
\item Nieto Gallegos Isaac Julián
\end{itemize}
\end{itemize}

Esta semana profundizamos en \textbf{administración de redes y conectividad remota} en Linux: SSH y su autenticación por llaves, el comando \texttt{ip} de \texttt{iproute2}, archivos de resolución de nombres y prácticas de confianza entre equipos en la red del laboratorio.
\subsection{Responsabilidades:}
\label{sec:org126e8c8}
\begin{itemize}
\item Gael Emiliano: SSH, Telnet y sección de máximas de seguridad.
\item Omar Fernando: Comando \texttt{ip}, rutas y vecinos; prácticas de red.
\item Mariana: \texttt{/etc/hosts}, \texttt{nsswitch.conf}, SELinux/semanage y estructura del documento.
\item Isaac Julián: Variables de entorno, loopback, actividades prácticas y observaciones.
\end{itemize}
\section{Conexión y acceso remoto}
\label{sec:org5a0f292}

\subsection{Conceptos nuevos}
\label{sec:orgfa0563d}

\subsubsection{Telnet}
\label{sec:org7d57aaa}
Telnet es un protocolo \textbf{histórico} (aprox. 1969, entorno ARPANET) que permite sesión remota en texto plano. Por convención usa el \textbf{puerto 23}. Hoy se considera \textbf{obsoleto para administración} porque \textbf{no cifra} usuario, contraseña ni datos: cualquier intermediario en la red puede leer la sesión.

\textbf{Uso razonable hoy:} diagnóstico puntual de conectividad (por ejemplo comprobar si un puerto responde o si hay ruta hasta un host), siempre entendiendo el riesgo. Para acceso remoto real en Linux se prefiere \textbf{SSH}.
\subsubsection{Autenticación SSH con llaves (\texttt{ssh-keygen})}
\label{sec:org477fdc7}
SSH (Secure Shell) es el estándar para acceso remoto \textbf{cifrado}; el servicio suele escuchar en el \textbf{puerto 22}. Además de contraseña, se puede usar un \textbf{par de llaves} (privada y pública):

\begin{verbatim}
ssh-keygen
\end{verbatim}

La llave pública (\texttt{id\_rsa.pub} u otro nombre elegido) se coloca en el servidor, en \texttt{\textasciitilde{}/.ssh/authorized\_keys} del usuario destino. Así el servidor puede verificar la identidad sin enviar la contraseña por la red, lo que \textbf{reduce ataques de fuerza bruta} y mejora la higiene operativa en laboratorios y servidores.
\subsection{Conceptos de repaso}
\label{sec:org9b64459}

\subsubsection{SSH}
\label{sec:orgc054038}
Protocolo de capa de aplicación para shell remoto, transferencia de archivos (\texttt{scp}, \texttt{sftp}) y túneles. El tráfico va cifrado; la identidad del servidor se valida típicamente contra \texttt{known\_hosts}.
\subsubsection{Ping}
\label{sec:org392ffce}
Herramienta ICMP (u otros modos según implementación) para probar alcanzabilidad y latencia. Complementa el trabajo con \texttt{/etc/hosts} cuando los nombres de compañeros ya están mapeados a IP.
\subsection{Máximas}
\label{sec:org4dce017}
\begin{itemize}
\item \textbf{Nunca administrar producción por Telnet}; usar SSH (y, cuando aplique, VPN).
\item \textbf{Proteger la llave privada} (\texttt{chmod 600} en \texttt{id\_rsa}, passphrase en entornos expuestos).
\item \textbf{Deshabilitar login root por contraseña} y preferir llaves o usuarios con \texttt{sudo}, según política del curso o del sitio.
\end{itemize}
\subsection{Sección libre}
\label{sec:orga8f8e92}
En laboratorio resulta útil combinar \textbf{ping por nombre} (vía \texttt{/etc/hosts}) con \textbf{SSH por IP o nombre} una vez que la resolución local es coherente. Si dos equipos se llaman igual en \texttt{hosts}, el primero que coincida en el archivo gana: conviene acordar nombres únicos por máquina.
\section{Administración de red en Linux: familia \texttt{ip} (\texttt{iproute2})}
\label{sec:org09cf9fe}

La suite \textbf{iproute2} reemplaza herramientas clásicas como \texttt{ifconfig} y \texttt{route} en muchas distribuciones modernas. El comando unificado es \texttt{ip}.
\subsection{Conceptos nuevos}
\label{sec:org12754dc}

\subsubsection{\texttt{ip addr show}}
\label{sec:orgedf5895}
Lista interfaces con direcciones \textbf{IPv4/IPv6} y, entre otros datos, la \textbf{MAC} asociada a cada interfaz.

\begin{verbatim}
ip addr show
\end{verbatim}
\subsubsection{\texttt{ip link}}
\label{sec:orgd8b8a1c}
Enfocado en la \textbf{capa de enlace}: estado \texttt{UP/DOWN}, tipo de interfaz, MTU, etc.

\begin{verbatim}
ip link
\end{verbatim}
\subsubsection{\texttt{ip route}}
\label{sec:orge9ded82}
Muestra la \textbf{tabla de enrutamiento} (gateway por defecto, rutas estáticas, alcances).

\begin{verbatim}
ip route
\end{verbatim}
\subsubsection{\texttt{ip neighbour}}
\label{sec:org3763a58}
Equivalente moderno de la tabla ARP: \textbf{direcciones MAC} de vecinos en segmentos locales alcanzables.

\begin{verbatim}
ip neighbour
\end{verbatim}
\subsection{Conceptos de repaso}
\label{sec:orgc10bcb0}
\begin{itemize}
\item Interfaz \textbf{loopback} (\texttt{lo}, típicamente \textbf{127.0.0.1}): tráfico del host consigo mismo; sirve para probar servicios locales (\texttt{sshd}, bases de datos, aplicaciones web en \texttt{localhost}) sin salir al cable.
\item Diferencia entre dirección \textbf{IP} (lógica, enrutable) y \textbf{MAC} (física/lógica de enlace en LAN).
\end{itemize}
\subsection{Máximas}
\label{sec:org4e32fdf}
\begin{itemize}
\item Antes de culpar a la aplicación, verificar \textbf{enlace} (\texttt{ip link}), \textbf{dirección} (\texttt{ip addr}) y \textbf{ruta} (\texttt{ip route}).
\end{itemize}
\subsection{Sección libre}
\label{sec:orgb19889f}
Si un comando antiguo aparece en un tutorial (\texttt{ifconfig}), en RHEL/Fedora/Debian reciente suele existir el equivalente con \texttt{ip}; instalar \texttt{net-tools} solo por nostalgia suele ser opcional.
\section{Configuración y seguridad del sistema}
\label{sec:org9d2c7c2}

\subsection{Conceptos nuevos}
\label{sec:orgbf25601}

\subsubsection{\texttt{/etc/nsswitch.conf}}
\label{sec:org98e8793}
Define el \textbf{orden} en que el sistema consulta fuentes para nombres de usuarios, grupos, hosts, etc. Para \textbf{hosts}, la línea relevante suele ser \texttt{hosts: files dns} (primero \texttt{/etc/hosts}, luego DNS). Cambiar ese orden altera el comportamiento de \texttt{ping}, \texttt{ssh} y cualquier programa que resuelva nombres.
\subsubsection{\texttt{semanage} (SELinux)}
\label{sec:orgf91d762}
En sistemas con \textbf{SELinux} en modo enforcing, \texttt{semanage} permite gestionar elementos de política: puertos, booleanos, contextos de archivos, etc. Es la capa correcta para \textbf{abrir un puerto} o ajustar reglas sin desactivar SELinux por completo (lo cual sería una mala práctica en entornos serios).
\subsection{Conceptos de repaso}
\label{sec:orgd537396}

\subsubsection{\texttt{/etc/hosts}}
\label{sec:orgcf4406a}
Archivo local que \textbf{mapea IP → nombre(s)}. Prioridad alta si \texttt{nsswitch} lo indica. Útil en aulas sin DNS interno o para nombres cortos entre compañeros (\texttt{ping compa1}, \texttt{ssh usuario@compa2}).
\subsection{Máximas}
\label{sec:org9523b9d}
\begin{itemize}
\item Coordinar con el grupo \textbf{qué IP corresponde a qué nombre} antes de llenar \texttt{hosts}; evita “funciona en mi máquina” por entradas inconsistentes.
\end{itemize}
\subsection{Sección libre}
\label{sec:org99087b3}

\subsubsection{Variables de entorno}
\label{sec:org1e17464}
Son pares \textbf{nombre=valor} heredados por procesos hijos. Afectan rutas de búsqueda de comandos (\texttt{PATH}), localización (\texttt{LANG}), opciones de \texttt{ssh} (\texttt{SSH\_AUTH\_SOCK}, etc.).

\begin{verbatim}
set | grep PATH
env | grep SSH
\end{verbatim}

Filtrar con \texttt{grep} ayuda cuando la salida de \texttt{set} es muy larga.
\section{Actividades prácticas reportadas}
\label{sec:org1d61359}

Resumen de lo trabajado en clase

\begin{enumerate}
\item \textbf{Mapeo de red:} Edición de \texttt{/etc/hosts} con las IPs de compañeros y nombres acordados para identificación en la red local.
\item \textbf{Intercambio de llaves:} Copia de \texttt{id\_rsa.pub} al \texttt{authorized\_keys} de otros nodos para establecer \textbf{confianza mutua} sin contraseña en cada inicio de sesión.
\item \textbf{Ejecución remota por SSH:} Una vez configurada la autenticación por llave, comandos de diagnóstico en equipos remotos, por ejemplo:
\end{enumerate}

\begin{verbatim}
ssh usuario@compañero 'uptime'
ssh usuario@compañero 'hostname'
ssh usuario@compañero 'who'
\end{verbatim}

\begin{enumerate}
\item \textbf{Análisis inalámbrico (conceptual):} Repaso de \textbf{ancho de banda} y de cómo el \textbf{SSID} identifica la red Wi-Fi en el punto de acceso; útil para relacionar la capa física/enlace con lo que vemos en \texttt{ip link} en interfaces \texttt{wlan*}.
\end{enumerate}
\end{document}
